\documentclass[11pt]{article}

\usepackage[margin=1in]{geometry}
\usepackage{amsmath, amssymb}
\usepackage{graphicx}
\usepackage{booktabs}
\usepackage[hidelinks]{hyperref}
\usepackage{placeins}

\title{BUSFIN 8200: Problem Set 2}
\author{Xi Wang}
\date{\today}

\begin{document}

\maketitle

\section*{Question 1}

\subsection*{(a)}

We have complete markets iff we can find a unique matrix $\Omega$ s.t.\ $X\Omega = I$, where
\[
X = \begin{bmatrix} 10 & 10 & 0 \\ 20 & 10 & 10 \\ 30 & 10 & 10 \end{bmatrix}, \qquad
\Omega = \begin{bmatrix} | & | & | \\ w_1 & w_2 & w_3 \\ | & | & | \end{bmatrix},
\]
and each $w_i$ is a portfolio weight vector using A, B and C. For example, $(10, 20, 30)'$ is asset A's payoff across the three states.

$X\Omega = I$ means
\[
X w_1 = \begin{pmatrix} 1 \\ 0 \\ 0 \end{pmatrix}, \qquad
X w_2 = \begin{pmatrix} 0 \\ 1 \\ 0 \end{pmatrix}, \qquad
X w_3 = \begin{pmatrix} 0 \\ 0 \\ 1 \end{pmatrix}.
\]

Because $X$ is $3 \times 3$, $X\Omega = I$ has a unique solution iff $X$ is invertible, and
\[
X \text{ invertible} \iff \det(X) \neq 0 \iff \operatorname{rank}(X) = 3.
\]
\[
X = \begin{bmatrix} 10 & 10 & 0 \\ 20 & 10 & 10 \\ 30 & 10 & 10 \end{bmatrix}
= 10 \begin{bmatrix} 1 & 1 & 0 \\ 2 & 1 & 1 \\ 3 & 1 & 1 \end{bmatrix}.
\]
Row reduction:
\begin{align*}
\text{row 3} - \text{row 2}: &\quad \begin{bmatrix} 1 & 1 & 0 \\ 1 & 0 & 0 \\ 3 & 1 & 1 \end{bmatrix} \\
\text{exchange rows 1 and 2}: &\quad \begin{bmatrix} 1 & 0 & 0 \\ 1 & 1 & 0 \\ 3 & 1 & 1 \end{bmatrix} \\
\text{row 2} - \text{row 1}: &\quad \begin{bmatrix} 1 & 0 & 0 \\ 0 & 1 & 0 \\ 3 & 1 & 1 \end{bmatrix} \\
\text{row 3} - 3\cdot\text{row 1} - \text{row 2}: &\quad \begin{bmatrix} 1 & 0 & 0 \\ 0 & 1 & 0 \\ 0 & 0 & 1 \end{bmatrix}
\end{align*}
Thus $\operatorname{rank}(X) = 3$ and $X$ is invertible. $\Omega = X^{-1}$ exists and is unique, with $X\Omega = XX^{-1} = I$. Thus, the market is complete.

\subsection*{(b)}

Let $X$ be the matrix of $N = 3$ non-redundant asset payoffs, let $P$ be the price vector of the original assets, and let $P_c = (q_1, q_2, q_3)'$ be the price vector of the contingent claims. Use the LOP to solve for $P_c$ as a function of $X$ and $P$:
\[
P = \begin{pmatrix} 16 \\ 8 \\ 6 \end{pmatrix}, \qquad
X = \begin{bmatrix} 10 & 10 & 0 \\ 20 & 10 & 10 \\ 30 & 10 & 10 \end{bmatrix}.
\]
The LOP says that two assets with identical payoffs in every state of nature must have the same price. Therefore,
\[
X' P_c = P \quad\Longrightarrow\quad
\begin{bmatrix} 10 & 20 & 30 \\ 10 & 10 & 10 \\ 0 & 10 & 10 \end{bmatrix}
\begin{pmatrix} q_1 \\ q_2 \\ q_3 \end{pmatrix}
= \begin{pmatrix} 16 \\ 8 \\ 6 \end{pmatrix}.
\]
Dividing everything by 10:
\[
\begin{cases}
q_1 + 2q_2 + 3q_3 = 1.6 \\
q_1 + q_2 + q_3 = 0.8 \\
q_2 + q_3 = 0.6
\end{cases}
\quad\Longrightarrow\quad
\begin{cases}
q_1 = 0.2 \\
q_2 = 0.4 \\
q_3 = 0.2
\end{cases}
\]
Thus, $P_c = (0.2, 0.4, 0.2)'$.

\subsection*{(c)}

\[
P_D = \text{Payoff}(D) \cdot P_c
= \begin{pmatrix} 30 & 20 & 10 \end{pmatrix}
\begin{pmatrix} 0.2 \\ 0.4 \\ 0.2 \end{pmatrix}
= 6 + 8 + 2 = 16.
\]

\subsection*{(d)}

\[
P(X) = E[MX] = \sum_{s=1}^{S} \pi(s) M(s) X(s).
\]
Since $q_s = \pi(s) M(s)$, we have $M(s) = q_s / \pi(s)$. With
\[
P_c = \begin{pmatrix} 0.2 \\ 0.4 \\ 0.2 \end{pmatrix}, \qquad
\pi(s) = \begin{pmatrix} 1/4 \\ 1/2 \\ 1/4 \end{pmatrix}
\quad\Longrightarrow\quad
M(s) = \begin{pmatrix} 0.8 \\ 0.8 \\ 0.8 \end{pmatrix}.
\]
By definition, no arbitrage means that
\begin{enumerate}
  \item An asset with payoff $\geq 0$ must have price $\geq 0$.
  \item An asset with payoff $\geq 0$ and a chance for payoff $> 0$ must have price $> 0$.
\end{enumerate}
We have $X \geq 0$, $\Pr(X > 0) > 0$, and $P(X) = E[MX] > 0$ (since $M > 0$). Thus, there is no arbitrage.

Moreover, $P(X) = E[MX] = 0.8\, E[X]$. The SDF stays constant across the three states, which implies it is risk-neutral.

\subsection*{(e)}

\begin{align*}
P_D = E[M X_D] &= \sum_{s=1}^{S} \pi(s) M(s) X_D(s) \\
&= \tfrac14 \cdot 0.8 \cdot 30 + \tfrac12 \cdot 0.8 \cdot 20 + \tfrac14 \cdot 0.8 \cdot 10 \\
&= 6 + 8 + 2 = 16.
\end{align*}
This is the same as in Q1(c).

\end{document}
